\section{Proving equivalence of X25519 in C and Coq}
\label{sec:C-Coq}

In this section we prove the following theorem:
% In this section we outline the structure of our proofs of the following theorem:

\begin{informaltheorem}
The implementation of X25519 in TweetNaCl (\TNaCle{crypto_scalarmult}) matches
the specifications of RFC~7748~\cite{rfc7748} (\Coqe{RFC}).
\end{informaltheorem}

More formally:
\begin{lstlisting}[language=Coq]
Theorem body_crypto_scalarmult:
  (* VST boiler plate. *)
  semax_body
    (* Clight translation of TweetNaCl. *)
    Vprog
    (* Hoare triples for fct calls. *)
    Gprog
    (* fct we verify. *)
    f_crypto_scalarmult_curve25519_tweet
    (* Our Hoare triple, see below. *)
    crypto_scalarmult_spec.
\end{lstlisting}

We first describe the global structure of our proof (\ref{subsec:proof-structure}).
Using our formalization of RFC~7748 (\sref{sec:Coq-RFC}) we specify the Hoare
triple before proving its correctness with VST (\ref{subsec:with-VST}).
We provide an example of equivalence of operations over different number
representations (\ref{subsec:num-repr-rfc}). Then, we describe efficient techniques
used to in some of our more complex proofs (\ref{subsec:inversions-reflections}).





\subsection{Structure of our proof}
\label{subsec:proof-structure}

In order to prove the correctness of X25519 in TweetNaCl code \TNaCle{crypto_scalarmult},
we use VST to prove that the code matches our functional Coq specification of \Coqe{RFC}.
Then, we prove that our specification of the scalar multiplication matches the mathematical definition
of elliptic curves and Theorem 2.1 by Bernstein~\cite{Ber06} (\sref{sec:maths}).

Verifying \TNaCle{crypto_scalarmult} also implies verifying all the functions
subsequently called: \TNaCle{unpack25519}; \TNaCle{A}; \TNaCle{Z}; \TNaCle{M};
\TNaCle{S}; \TNaCle{car25519}; \TNaCle{inv25519}; \TNaCle{set25519}; \TNaCle{sel25519};
\TNaCle{pack25519}.

We prove the implementation of X25519 is \textbf{sound}, \ie:
\begin{itemize}
\item absence of access out-of-bounds of arrays (memory safety).
\item absence of overflows/underflow in the arithmetic.
\end{itemize}
We also prove that TweetNaCl's code is \textbf{correct}:
\begin{itemize}
\item X25519 is correctly implemented (we get what we expect) .
\item Operations on \TNaCle{gf} (\TNaCle{A}, \TNaCle{Z}, \TNaCle{M}, \TNaCle{S})
are equivalent to operations ($+,-,\times,x^2$) in $\Zfield$.
\item The Montgomery ladder computes the multiple of a point.
% \item The Montgomery ladder computes a scalar multiplication between a natural
% number and a point.
\end{itemize}

In order to prove the soundness and correctness of \TNaCle{crypto_scalarmult},
we reuse the generic Montgomery ladder defined in \sref{sec:Coq-RFC}.

We define a High-level specification by instantiating the ladder with a generic
field $\K$, this allows us to prove the correctness of the ladder with respect
to the theory of elliptic curves.
This high-level specification does not rely on the parameters of Curve25519.
We later specialize $\K$ with $\Ffield$, and the parameters of Curve25519 ($a = 486662, b = 1$),
to derive the correctness of \coqe{RFC} (\sref{sec:maths}).

We define a mid-level specification by instantiating the ladder over $\Zfield$.
Additionally we also provide a low-level specification close to the \texttt{C} code
(over lists of $\Z$). We show this specification to be equivalent to the
\emph{semantic version} of C (Clight) using the VST.
This low level specification gives us the soundness assurance.

RFC~7748's X25519 formalization (\sref{sec:Coq-RFC}) takes as input list of $\Z$.
However the inner Montgomery ladder operates on $\Zfield$. We show its equivalence
with our mid-level and low-level specifications.

By showing that operations over instances ($\K = \Ffield$, $\Zfield$, list of $\Z$) are
equivalent, we bridge the gap between the different level of specification
with Curve25519 parameters.
As such, we prove all specifications to equivalent (\fref{tikz:ProofStructure}).
This guarantees us the correctness of the implementation.

\begin{figure}[h]
  \centering
  \begin{tikzpicture}[textstyle/.style={black, anchor= south west, align=center}]

    \filldraw[draw=orange!10!doc@lstbackground, fill=doc@lstbackground, thick] (0.25,0.5) rectangle (4.5,5.5);
    % node[textstyle, anchor=west, draw=yellow, fill=yellow!20, thick, minimum width=5.5cm,minimum height=5cm] {};

    \draw (4.5,5.5)  node[anchor=north east, inner sep=0pt] (russell) {\includegraphics[width=.03\textwidth]{img/coq_logo.png}};

    \draw (0.5,-1) node [textstyle, anchor=west, draw=black, thick, minimum width=3cm,minimum height=0.5cm] (longlong) {\texttt{long long[16]}};
    \draw (0,-1) node [anchor=east] (longlongdef) {\texttt{C} code};
    \draw (2.25,-0.1) node [anchor=west] (app) {\texttt{clightgen}};


    \draw (0.5,1) node [textstyle, anchor=west, draw=black, thick, minimum width=3cm,minimum height=0.5cm] (clight) {\texttt{tptr tlong}};
    \draw (0,1) node [anchor=east] (clightdef) {Clight};

    \draw [thick, ->] (longlong.north) -- (clight.south);

    \begin{scope}[yshift=1 cm,xshift=0 cm]

      \draw (0.5,2) node [textstyle, anchor=west, draw=black, thick, minimum width=3cm,minimum height=0.5cm] (ll) {\texttt{list} $\Z$};
      \draw (0,2) node [anchor=east] (shn) {Low Level};

      \draw [thick,double, <->, >=implies] (clight.north) -- (ll.south);
      \draw (2.25,1) node[anchor=west, inner sep=0pt] (chain) {\includegraphics[width=.07\textwidth]{img/chain.png}};


      \draw (0.5,3) node [textstyle, anchor=west, draw=black, thick, minimum width=3cm,minimum height=0.5cm] (ml) {$\Zfield$};
      \draw (0,3) node [anchor=east] (shn) {Mid Level};

      \draw[thick,double, <->, >=implies] (ll.north) -- (ml.south);

      \draw (0.5,4) node [textstyle, anchor=west, draw=black, thick, minimum width=3cm,minimum height=0.5cm] (hl) {$\K$};
      \draw (0,4) node [anchor=east] (shn) {High Level};

      \draw[thick,double, <-, >=implies] (ml.north) -- (hl.south);
    \end{scope}
\end{tikzpicture}

  \caption{Structural construction of the proof}
  \label{tikz:ProofStructure}
\end{figure}






% The location of the definitions are listed in Appendix~\ref{appendix:proof-files}.






\subsection{Applying the Verifiable Software Toolchain}
\label{subsec:with-VST}

We now turn our focus to the formal specification of \TNaCle{crypto_scalarmult}.
We use our definition of X25519 from the RFC in the Hoare triple and prove
its correctness.

\subheading{Specifications.}
We show the soundness of TweetNaCl by proving the following specification matches
a pure Coq function.
% why "pure" ?
% A pure function is a function where the return value is only determined by its
% input values, without observable side effects (Side effect are e.g. printing)
This defines the equivalence between the Clight representation and our Coq
definition of the ladder (\coqe{RFC}).

\begin{lstlisting}[language=CoqVST]
Definition crypto_scalarmult_spec :=
DECLARE _crypto_scalarmult_curve25519_tweet
WITH
  v_q: val, v_n: val, v_p: val, c121665:val,
  sh : share,
  q : list val, n : list Z, p : list Z
(*------------------------------------------*)
PRE [ _q OF (tptr tuchar),
     _n OF (tptr tuchar),
     _p OF (tptr tuchar) ]
PROP (writable_share sh;
      Forall (fun x => 0 <= x < 2^8) p;
      Forall (fun x => 0 <= x < 2^8) n;
      Zlength q = 32; Zlength n = 32;
      Zlength p = 32)
LOCAL(temp _q v_q; temp _n v_n; temp _p v_p;
      gvar __121665 c121665)
SEP  (sh [{ v_q }] <<(uch32)-- q;
      sh [{ v_n }] <<(uch32)-- mVI n;
      sh [{ v_p }] <<(uch32)-- mVI p;
      Ews [{ c121665 }] <<(lg16)-- mVI64 c_121665)
(*------------------------------------------*)
POST [ tint ]
PROP (Forall (fun x => 0 <= x < 2^8) (RFC n p);
      Zlength (RFC n p) = 32)
LOCAL(temp ret_temp (Vint Int.zero))
SEP  (sh [{ v_q }] <<(uch32)-- mVI (RFC n p);
      sh [{ v_n }] <<(uch32)-- mVI n;
      sh [{ v_p }] <<(uch32)-- mVI p;
      Ews [{ c121665 }] <<(lg16)-- mVI64 c_121665
\end{lstlisting}

In this specification we state as preconditions:
\begin{itemize}
  \item[] \VSTe{PRE}: \VSTe{_p OF (tptr tuchar)}\\
  The function \TNaCle{crypto_scalarmult} takes as input three pointers to
  arrays of unsigned bytes (\VSTe{tptr tuchar}) \VSTe{_p}, \VSTe{_q} and \VSTe{_n}.
  \item[] \VSTe{LOCAL}: \VSTe{temp _p v_p}\\
  Each pointer represent an address \VSTe{v_p},
  \VSTe{v_q} and \VSTe{v_n}.
  \item[] \VSTe{SEP}: \VSTe{sh [{ v_p $\!\!\}\!\!]\!\!\!$ <<(uch32)-- mVI p}\\
  In the memory share \texttt{sh}, the address \VSTe{v_p} points
  to a list of integer values \VSTe{mVI p}.
  \item[] \VSTe{PROP}: \VSTe{Forall (fun x => 0 <= x < 2^8) p}\\
  In order to consider all the possible inputs, we assume each
  elements of the list \texttt{p} to be bounded by $0$ included and $2^8$
  excluded.
  \item[] \VSTe{PROP}: \VSTe{Zlength p = 32}\\
  We also assume that the length of the list \texttt{p} is 32. This defines the
  complete representation of \TNaCle{u8[32]}.
\end{itemize}

As Post-condition we have:
\begin{itemize}
  \item[] \VSTe{POST}: \VSTe{tint}\\
  The function \TNaCle{crypto_scalarmult} returns an integer.
  \item[] \VSTe{LOCAL}: \VSTe{temp ret_temp (Vint Int.zero)}\\
  The returned integer has value $0$.
  \item[] \VSTe{SEP}: \VSTe{sh [{ v_q $\!\!\}\!\!]\!\!\!$ <<(uch32)-- mVI (RFC n p)}\\
  In the memory share \texttt{sh}, the address \VSTe{v_q} points
  to a list of integer values \VSTe{mVI (RFC n p)} where \VSTe{RFC n p} is the
  result of the \TNaCle{crypto_scalarmult} of \VSTe{n} and \VSTe{p}.
  \item[] \VSTe{PROP}: \VSTe{Forall (fun x => 0 <= x < 2^8) (RFC n p)}\\
  \VSTe{PROP}: \VSTe{Zlength (RFC n p) = 32}\\
  We show that the computation for \VSTe{RFC} fits in  \TNaCle{u8[32]}.
\end{itemize}

This specification (proven with VST) shows that \TNaCle{crypto_scalarmult} in C
computes the same result as \VSTe{RFC} in Coq provided that inputs are within
their respective bounds: arrays of 32 bytes.

The correctness of this specification is formally proven in Coq with
\coqe{Theorem body_crypto_scalarmult}.

% \begin{theorem}
% \label{thm:crypto-vst}
% \TNaCle{crypto_scalarmult} in TweetNaCl has the same behavior as \coqe{RFC} in Coq.
% \end{theorem}



% % this does not exists anymore
% The location of the proof of this Hoare triple can be found in Appendix~\ref{appendix:proof-files}.

% We now bring the attention of the reader to details of verifications using the VST.

%% BLABLA about VST. Does not belong here.

% The Verifiable Software Toolchain uses a strongest postcondition strategy.
% The user must first write a formal specification of the function he wants to verify in Coq.
% This should be as close as possible to the C implementation behavior.
% This will simplify the proof and help with stepping through the Clight version of the software.
% With the range of inputs defined, VST mechanically steps through each instruction
% and ask the user to verify auxiliary goals such as array bound access, or absence of overflows/underflows.
% We call this specification a low level specification. A user will then have an easier
% time to prove that his low level specification matches a simpler higher level one.

% In order to further speed-up the verification process, it has to be know that to
% prove the specification \TNaCle{crypto_scalarmult}, a user only need the specification of e.g. \TNaCle{M}.
% This provide with multiple advantages: the verification by the Coq kernel can be done
% in parallel and multiple users can work on proving different functions at the same time.
% For the sake of completeness we proved all intermediate functions.

\subheading{Memory aliasing.}
In the VST, a simple specification of \texttt{M(o,a,b)} will assume three
distinct memory share. When called with three memory shares (\texttt{o, a, b}),
the three of them will be consumed. However assuming this naive specification
when \texttt{M(o,a,a)} is called (squaring), the first two memory shares (\texttt{o, a})
are consumed and VST will expect a third memory share (\texttt{a}) which does not \emph{exist} anymore.
Examples of such cases are illustrated in \fref{tikz:MemSame}.
\begin{figure}[h]
  \centering
  \begin{tikzpicture}
\draw (0,0) rectangle (3,0.3);
\draw[thick] (1,0) -- (1,0.3);
\draw[thick] (2,0) -- (2,0.3);
\node [anchor=east] (shn) at (0,0.15) {sh};
\node [anchor=north] (Moab) at (1.5,-0.5) {\texttt{M(o,a,b)}};
\draw[thick, -> ] (1.25, -0.6) -- (1.25, -0.5) -- (0.5,-0.1);
\draw[thick, -> ] (1.6, -0.6) -- (1.6, -0.5) -- (1.5,-0.1);
\draw[thick, -> ] (1.95, -0.6) -- (1.95, -0.5) -- (2.5,-0.1);

\node [anchor=north west, text width=6cm, align=left] (Sepoab) at (-0.5,-1.8)
  {\VSTess{sh [{ v_o $\!\!\}\!\!]\!\!\!$ <<(lg16)-- o}\\[-0.2em]
  \VSTess{sh [{ v_a $\!\!\}\!\!]\!\!\!$ <<(lg16)-- a}\\[-0.2em]
  \VSTess{sh [{ v_b $\!\!\}\!\!]\!\!\!$ <<(lg16)-- b}};


\begin{scope}[yshift=0 cm,xshift=4.5 cm]
  \draw (0,0) rectangle (2,0.3);
  \draw[thick] (1,0) -- (1,0.3);
  % \draw[thick] (2,0) -- (2,0.3);
  \node [anchor=east] (shm) at (0,0.15) {sh};
  \node [anchor=north] (Mcaa) at (1,-0.5) {\texttt{M(c,c,a)}};
  \draw[thick, -> ] (0.75, -0.6) -- (0.75, -0.5) -- (0.4,-0.1);
  \draw[thick, -> ] (1.1, -0.6) -- (1.1, -0.5) -- (0.6,-0.1);
  \draw[thick, -> ] (1.45, -0.6) -- (1.45, -0.5) -- (1.5,-0.1);

  \node [anchor=north west, text width=6cm, align=left] (Sepoab) at (-0.5,-1.8)
    {\VSTess{sh [{ v_c $\!\!\}\!\!]\!\!\!$ <<(lg16)-- c}\\[-0.2em]
    \VSTess{sh [{ v_a $\!\!\}\!\!]\!\!\!$ <<(lg16)-- a}};
\end{scope}

\draw[dashed] (-0.5,-1.2) -- +(8,0);
\node [anchor=north west] (sep) at (-0.5,-1.2) {In Separation Logic:};

\begin{scope}[yshift=-4.25 cm,xshift=2 cm]
  \draw (0,0) rectangle (2,0.3);
  \draw[thick] (1,0) -- (1,0.3);
  \node [anchor=east] (shm) at (0,0.15) {sh\;};

  \draw (0,0.5) rectangle (1,0.8);
  \node [anchor=east] (shm) at (0,0.65) {Ews};

  \draw[thick] (2,0) -- (2,0.3);
  \node [anchor=north west] (Mcaa) at (0.15,-0.5) {\texttt{M(a,c,\_121665)}};
  \draw[thick, -> ] (0.75, -0.6) -- (0.75, -0.5) -- (0.5,-0.1);
  \draw[thick, -> ] (1.1, -0.6) -- (1.1, -0.5) -- (1.5,-0.1);
  \draw[thick, -> ] (2.0, -0.6) -- (2.45, -0.2) -- (2.45, 0.3) -- (2.2, 0.65) -- (1.5,0.65);

  \node [anchor=north west, text width=6cm, align=left] (Sepoab) at (-0.5,-1.8)
    {\VSTess{sh [{ v_a $\!\!\}\!\!]\!\!\!$ <<(lg16)-- a}\\[-0.2em]
    \VSTess{sh [{ v_c $\!\!\}\!\!]\!\!\!$ <<(lg16)-- c}\\[-0.2em]
    \VSTess{Ews [{ __121665 $\!\!\}\!\!]\!\!\!$ <<(lg16)-- _121665}
    };
\end{scope}
\draw[dashed] (0.5,-5.45) -- +(6,0);
\node [anchor=north west] (sep) at (0.5,-5.45) {In Separation Logic:};
\end{tikzpicture}%

  \caption{Aliasing and Separation Logic}
  \label{tikz:MemSame}
\end{figure}

As a result, a function must either have multiple specifications or specify which
aliasing case is being used.
The first option would require us to do very similar proofs multiple times for a same function.
We chose the second approach: for functions with 3 arguments, named hereafter \texttt{o, a, b},
we define an additional parameter $k$ with values in $\{0,1,2,3\}$:
\begin{itemize}
  \item if $k=0$ then \texttt{o} and \texttt{a} are aliased.
  \item if $k=1$ then \texttt{o} and \texttt{b} are aliased.
  \item if $k=2$ then \texttt{a} and \texttt{b} are aliased.
  \item else there is no aliasing.
\end{itemize}
In the proof of our specification, we do a case analysis over $k$ when needed.
This solution does not cover all cases (e.g. all arguments are aliased) but it
is enough for our needs.

\subheading{Verifying \texttt{for} loops.}
Final states of \texttt{for} loops are usually computed by simple recursive functions.
However we must define invariants which are true for each iteration step.

Assume we want to prove a decreasing loop where indexes go from 3 to 0.
Define a function $g : \N \rightarrow State  \rightarrow State $ which takes as
input an integer for the index and a state, then returns a state.
It simulates the body of the \texttt{for} loop.
Assume its recursive call: $f : \N \rightarrow State \rightarrow State $ which
iteratively applies $g$ with decreasing index:
\begin{equation*}
  f ( i , s ) =
  \begin{cases}
  s & \text{if } s = 0 \\
  f( i - 1 , g ( i - 1  , s )) & \text{otherwise}
  \end{cases}
\end{equation*}
Then we have :
\begin{align*}
  f(4,s) &= g(0,g(1,g(2,g(3,s))))
  % \\
  % f(3,s) &= g(0,g(1,g(2,s)))
\end{align*}
To prove the correctness of $f(4,s)$, we need to prove that intermediate steps
$g(3,s)$; $g(2,g(3,s))$; $g(1,g(2,g(3,s)))$; $g(0,g(1,g(2,g(3,s))))$ are correct.
Due to the computation order of recursive function, our loop invariant for
$i\in\{0;1;2;3;4\}$ cannot use $f(i)$.
To solve this, we define an auxiliary function with an accumulator such that
given $i\in\{0;1;2;3;4\}$, it will compute the first $i$ steps of the loop.

We then prove for the complete number of steps, the function with the accumulator
and without returns the same result.
We formalized this result in a generic way in Appendix~\ref{subsubsec:for}.

Using this formalization, we prove that the 255 steps of the Montgomery ladder
in C provide the same computations as in \coqe{RFC}.





\subsection{Number Representation and C Implementation}
\label{subsec:num-repr-rfc}

As described in \sref{subsec:Number-TweetNaCl}, numbers in \TNaCle{gf} are represented
in base $2^{16}$ and we use a direct mapping to represent that array as a list
integers in Coq. However in order to show the correctness of the basic operations,
we need to convert this number to a full integer.
We reuse the mapping
$\text{\coqe{ZofList}} : \Z \rightarrow \texttt{list}~\Z \rightarrow \Z$ from \sref{subsec:integer-bytes}
and define a notation where $n$ is $16$, placing us with a radix of $2^{16}$.
\begin{lstlisting}[language=Coq]
Notation "Z16.lst A" := (ZofList 16 A).
\end{lstlisting}
To facilitate working in $\Zfield$, we define the \coqe{:GF} notation.
\begin{lstlisting}[language=Coq]
Notation "A :GF" := (A mod (2^255-19)).
\end{lstlisting}
Later in \sref{subsec:Zmodp}, we formally define $\Ffield$.
Equivalence between operations in $\Zfield$ (\ie under \coqe{:GF}) and in $\Ffield$ are easily proven.

Using these two definitions, we prove intermediate lemmas such as the correctness of the
multiplication \Coqe{Low.M} where \Coqe{Low.M} replicate the computations and steps done in C.
\begin{lemma}
  \label{lemma:mult_correct}
  \Coqe{Low.M} implements correctly the multiplication over $\Zfield$.
\end{lemma}
And specified in Coq as follows:
\begin{lstlisting}[language=Coq]
Lemma mult_GF_Zlength :
  forall (a:list Z) (b:list Z),
  Zlength a = 16 ->
  Zlength b = 16 ->
   (Z16.lst (Low.M a b)) :GF =
   (Z16.lst a * Z16.lst b) :GF.
\end{lstlisting}

However for our purpose, simple functional correctness is not enough.
We also need to define the bounds under which the operation is correct,
allowing us to chain them, guaranteeing us the absence of overflow.

\begin{lemma}
  \label{lemma:mult_bounded}
  if all the values of the input arrays are constrained between $-2^{26}$ and $2^{26}$,
  then the output of \coqe{Low.M} will be constrained between $-38$ and $2^{16} + 38$.
\end{lemma}
And seen in Coq as follows:
\begin{lstlisting}[language=Coq]
Lemma M_bound_Zlength :
  forall (a:list Z) (b:list Z),
  Zlength a = 16 ->
  Zlength b = 16 ->
  Forall (fun x => -2^26 < x < 2^26) a ->
  Forall (fun x => -2^26 < x < 2^26) b ->
    Forall (fun x => -38 <= x < 2^16 + 38) (Low.M a b).
\end{lstlisting}





\subsection{Inversions, Reflections and Packing}
\label{subsec:inversions-reflections}

We now turn our attention to the multiplicative inverse in $\Zfield$ and techniques
to improve the verification speed of complex formulas.

\subheading{Inversion in \Zfield.}
We define a Coq version of the inversion mimicking
the behavior of \TNaCle{inv25519} (see below) over \Coqe{list Z}.
\begin{lstlisting}[language=Ctweetnacl]
sv inv25519(gf o,const gf i)
{
  gf c;
  int a;
  set25519(c,i);
  for(a=253;a>=0;a--) {
    S(c,c);
    if(a!=2&&a!=4) M(c,c,i);
  }
  FOR(a,16) o[a]=c[a];
}
\end{lstlisting}
We specify this with 2 functions: a recursive \Coqe{pow_fn_rev}
to simulate the \texttt{for} loop and a simple \Coqe{step_pow} containing the body.
\begin{lstlisting}[language=Coq]
Definition step_pow (a:Z)
  (c:list Z) (g:list Z) : list Z :=
  let c := Sq c in
    if a <>? 2 && a <>? 4
    then M c g
    else c.

Function pow_fn_rev (a:Z) (b:Z)
  (c: list Z) (g: list Z)
  {measure Z.to_nat a} : (list Z) :=
  if a <=? 0
    then c
    else
      let prev := pow_fn_rev (a - 1) b c g in
        step_pow (b - a) prev g.
\end{lstlisting}
This \Coqe{Function} requires a proof of termination. It is done by proving the
well-foundedness of the decreasing argument: \Coqe{measure Z.to_nat a}. Calling
\Coqe{pow_fn_rev} 254 times allows us to reproduce the same behavior as the Clight definition.
\begin{lstlisting}[language=Coq]
Definition Inv25519 (x:list Z) : list Z :=
  pow_fn_rev 254 254 x x.
\end{lstlisting}
Similarly we define the same function over $\Z$.
\begin{lstlisting}[language=Coq]
Definition step_pow_Z (a:Z) (c:Z) (g:Z) : Z :=
  let c := c * c in
  if a <>? 2 && a <>? 4
    then c * g
    else c.

Function pow_fn_rev_Z (a:Z) (b:Z) (c:Z) (g: Z)
  {measure Z.to_nat a} : Z :=
  if (a <=? 0)
    then c
    else
      let prev := pow_fn_rev_Z (a - 1) b c g in
        step_pow_Z (b - a) prev g.

Definition Inv25519_Z (x:Z) : Z :=
  pow_fn_rev_Z 254 254 x x.
\end{lstlisting}
By using \lref{lemma:mult_correct}, we prove their equivalence in $\Zfield$.
\begin{lemma}
\label{lemma:Inv_equivalence}
The function \coqe{Inv25519} over list of integers computes the same
result at \coqe{Inv25519_Z} over integers in \Zfield.
\end{lemma}
This is formalized in Coq as follows:
\begin{lstlisting}[language=Coq]
Lemma Inv25519_Z_GF : forall (g:list Z),
  length g = 16 ->
  (Z16.lst (Inv25519 g)) :GF =
  (Inv25519_Z (Z16.lst g)) :GF.
\end{lstlisting}

In TweetNaCl, \TNaCle{inv25519} computes an inverse in $\Zfield$.
It uses Fermat's little theorem by doing an exponentiation to $2^{255}-21$.
This is done by applying a square-and-multiply algorithm. The binary representation
of $p-2$ implies to always do multiplications except for bits 2 and 4.
To prove the correctness of the result we can use multiple strategies such as:
\begin{itemize}
  \item Proving it is a special case of square-and-multiply algorithm applied to $2^{255}-21$.
  \item Unrolling the for loop step-by-step and applying the equalities
  $x^a \times x^b = x^{(a+b)}$ and $(x^a)^2 = x^{(2 \times a)}$. We prove that
  the resulting exponent is $2^{255}-21$.
\end{itemize}
We use the second method for the benefits of simplicity. However it requires to
apply the unrolling and exponentiation formulas 255 times. This can be automated
in Coq with tacticals such as \Coqe{repeat}, but it generates a proof object which
will take a long time to verify.

\subheading{Speeding up with Reflections.} This technique provides us with
flexibility, \eg we don't need to know the number of times nor the order
in which the lemmas needs to be applied (chapter 15 in \cite{CpdtJFR}).

The idea is to \emph{reflect} the goal into a decidable environment.
We show that for a property $P$, we can define a decidable Boolean property
$P_{bool}$ such that if $P_{bool}$ is \Coqe{true} then $P$ holds.
$$reify\_P : P_{bool} = true \implies P$$
By applying $reify\_P$ on $P$ our goal become $P_{bool} = true$.
We then compute the result of $P_{bool}$. If the decision goes well we are
left with the tautology $true = true$.

To prove that the \Coqe{Inv25519_Z} is computing $x^{2^{255}-21}$,
we define a Domain Specific Language (DSL).
\begin{dfn}
Let \Coqe{expr_inv} denote an expression which is either a term;
a multiplication of expressions; a squaring of an expression or a power of an expression.
And Let \Coqe{formula_inv} denote an equality between two expressions.
\end{dfn}
\begin{lstlisting}[language=Coq]
Inductive expr_inv :=
  | R_inv : expr_inv
  | M_inv : expr_inv -> expr_inv -> expr_inv
  | S_inv : expr_inv -> expr_inv
  | P_inv : expr_inv -> positive -> expr_inv.

Inductive formula_inv :=
  | Eq_inv : expr_inv -> expr_inv -> formula_inv.
\end{lstlisting}
The denote functions are defined as follows:
\begin{lstlisting}[language=Coq]
Fixpoint e_inv_denote (m:expr_inv) : Z :=
  match m with
  | R_inv     =>
    term_denote
  | M_inv x y =>
    (e_inv_denote x) * (e_inv_denote y)
  | S_inv x =>
    (e_inv_denote x) * (e_inv_denote x)
  | P_inv x p =>
    pow (e_inv_denote x) (Z.pos p)
  end.

Definition f_inv_denote (t : formula_inv) : Prop :=
  match t with
  | Eq_inv x y => e_inv_denote x = e_inv_denote y
  end.
\end{lstlisting}
All denote functions also take as an argument the environment containing the variables.
We do not show it here for the sake of readability,
given that an environment, \Coqe{term_denote} returns the appropriate variable.
With this Domain Specific Language we have the following equality:
\begin{lstlisting}[backgroundcolor=\color{white}]
f_inv_denote
 (Eq_inv (M_inv R_inv (S_inv R_inv))
         (P_inv R_inv 3))
  = (x * x^2 = x^3)
\end{lstlisting}
On the right side, \Coqe{(x * x^2 = x^3)} depends on $x$. On the left side,
\texttt{(Eq\_inv (M\_inv R\_inv (S\_inv R\_inv)) (P\_inv R\_inv 3))} does not depend on $x$.
This allows us to use computations in our decision procedure.

We define \Coqe{step_inv} and \Coqe{pow_inv} to mirror the behavior of
\Coqe{step_pow_Z} and respectively \Coqe{pow_fn_rev_Z} over our DSL and
we prove their equality.
\begin{lstlisting}[language=Coq]
Lemma step_inv_step_pow_eq :
  forall (a:Z) (c:expr_inv) (g:expr_inv),
  e_inv_denote (step_inv a c g) =
  step_pow_Z a (e_inv_denote c) (e_inv_denote g).

Lemma pow_inv_pow_fn_rev_eq :
  forall (a:Z) (b:Z) (c:expr_inv) (g:expr_inv),
  e_inv_denote (pow_inv a b c g) =
  pow_fn_rev_Z a b (e_inv_denote c) (e_inv_denote g).
\end{lstlisting}
We then derive the following lemma.
\begin{lemma}
\label{lemma:reify}
With an appropriate choice of variables, \Coqe{pow_inv} denotes \Coqe{Inv25519_Z}.
\end{lemma}

In order to prove formulas in \Coqe{formula_inv},
we have the following a decidable procedure.
We define \Coqe{pow_expr_inv}, a function which returns the power of an expression.
We then compare the two values and decide over their equality.
\begin{lstlisting}[language=Coq]
Fixpoint pow_expr_inv (t:expr_inv) : Z :=
  match t with
  (* power of a term is 1. *)
  | R_inv   => 1
  (* x^a * x^b = x^{a+b}. *)
  | M_inv x y =>
    (pow_expr_inv x) + (pow_expr_inv y)
    (* (x^a)^2 = x^{2a}. *)
  | S_inv x =>
    2 * (pow_expr_inv x)
    (* (x^b)^a = x^{a*b}. *)
  | P_inv x p =>
    (Z.pos p) * (pow_expr_inv x)
  end.

Definition decide_e_inv (l1 l2:expr_inv) : bool :=
  (pow_expr_inv l1) ==? (pow_expr_inv l2).

Definition decide_f_inv (f:formula_inv) : bool :=
  match f with
  | Eq_inv x y => decide_e_inv x y
  end.
\end{lstlisting}

We prove our decision procedure correct.
\begin{lemma}
\label{lemma:decide}
For all formulas $f$, if the decision over $f$ returns \Coqe{true},
then the denoted equality by $f$ is true.
\end{lemma}
Which is formalized as:
\begin{lstlisting}[language=Coq]
Lemma decide_formula_inv_impl :
  forall (f:formula_inv),
  decide_f_inv f = true ->
    f_inv_denote f.
\end{lstlisting}

By reification to over DSL (\lref{lemma:reify}) and by applying our decision
(\lref{lemma:decide}), we prove the following corollary.
\begin{lemma}
\label{lemma:inv_comput_inv}
\Coqe{Inv25519_Z} computes an inverse in \Zfield.
\end{lemma}
Which is formalized as:
\begin{lstlisting}[language=Coq]
Theorem Inv25519_Z_correct :
  forall (x:Z),
  Inv25519_Z x = pow x (2^255-21).
\end{lstlisting}

% From \Coqe{Inv25519_Z_GF} (\lref{lemma:Inv_equivalence}) and \Coqe{Inv25519_Z_correct} (Corollary~\ref{lemma:inv_comput_inv}),
From \lref{lemma:Inv_equivalence} and Corollary~\ref{lemma:inv_comput_inv},
we conclude the functional correctness of the inversion over \Zfield.
\begin{corollary}
\label{cor:inv_comput_field}
\Coqe{Inv25519} computes an inverse in \Zfield.
\end{corollary}
Which is formalized as:
\begin{lstlisting}[language=Coq]
Corollary Inv25519_Zpow_GF :
  forall (g:list Z),
  length g = 16 ->
  Z16.lst (Inv25519 g) :GF  =
  (pow (Z16.lst g) (2^255-21)) :GF.
\end{lstlisting}

\subheading{Another applications of reflections: packing}
This reflection technique can also be used where proofs requires some computing
over a small and finite domain of variables to test e.g. for all $i$ such that $0 \le i < 16$.
Using reflection we prove that we can split the for loop in \TNaCle{pack25519} into two parts.
\begin{lstlisting}[language=Ctweetnacl]
for(i=1;i<15;i++) {
  m[i]=t[i]-0xffff-((m[i-1]>>16)&1);
  m[i-1]&=0xffff;
}
\end{lstlisting}
The first loop is computing the subtraction while the second is applying the carries.
\begin{lstlisting}[language=Ctweetnacl]
for(i=1;i<15;i++) {
  m[i]=t[i]-0xffff
}

for(i=1;i<15;i++) {
  m[i]=m[i]-((m[i-1]>>16)&1);
  m[i-1]&=0xffff;
}
\end{lstlisting}

This loop separation allows simpler proofs. The first loop is seen as the subtraction of a number in \Zfield.
We then prove that with the iteration of the second loop, the number represented in $\Zfield$ stays the same.
This leads to the proof that \TNaCle{pack25519} is effectively reducing modulo $\p$ and returning a number in base $2^8$.
\begin{lstlisting}[language=Coq]
Lemma Pack25519_mod_25519 :
  forall (l:list Z),
  Zlength l = 16 ->
  Forall (fun x => -2^62 < x < 2^62) l ->
  ZofList 8 (Pack25519 l) =
  (Z16.lst l) mod (2^255-19).
\end{lstlisting}

By using each functions \coqe{Low.M}; \coqe{Low.A}; \coqe{Low.Sq}; \coqe{Low.Zub};
\coqe{Unpack25519}; \coqe{clamp}; \coqe{Pack25519}; \coqe{Inv25519}; \coqe{car25519}; \coqe{montgomery_rec},
we defined a Coq definition \coqe{Crypto_Scalarmult} mimicking the exact behavior of X25519 in TweetNaCl.

By proving that each functions \coqe{Low.M}; \coqe{Low.A}; \coqe{Low.Sq}; \coqe{Low.Zub};
\coqe{Unpack25519}; \coqe{clamp}; \coqe{Pack25519}; \coqe{Inv25519}; \coqe{car25519} are behaving over \coqe{list Z}
as their equivalent over \coqe{Z} in \coqe{:GF} (in \Zfield), we prove that given the same inputs \coqe{Crypto_Scalarmult} applies the same computation as \coqe{RFC}.
% This is formalized as follows in Coq:
\begin{lstlisting}[language=Coq]
Lemma Crypto_Scalarmult_RFC_eq :
  forall (n: list Z) (p: list Z),
  Zlength n = 32 ->
  Zlength p = 32 ->
  Forall (fun x => 0 <= x /\ x < 2 ^ 8) n ->
  Forall (fun x => 0 <= x /\ x < 2 ^ 8) p ->
  Crypto_Scalarmult n p = RFC n p.
\end{lstlisting}

This proves that TweetNaCl's X25519 implementation respect RFC~7748.
